% ============================================================================
\chapter{Results of Morphological Analysis}
\label{ch:morphology}
% ============================================================================

This chapter reports the results of the paired three-dimensional morphological analysis of the 45 tracer particles---five per configuration across the nine configurations of the experimental programme. The chapter is organised in three sections. \Cref{sec:volume-change} reports the per-particle volume-change distribution and identifies a small set of outlier particles using the Tukey inter-quartile-range criterion described in \cref{sec:tukey}. \Cref{sec:shape-evolution} reports the paired evolution of four principal shape descriptors---sphericity, convexity, angularity, and roundness---before and after the impact loading. \Cref{sec:morph-discussion} examines the global pattern of mean change across all twenty-two morphological descriptors using two heat maps.

\section{Per-particle volume change and outlier identification}
\label{sec:volume-change}

Across the 45 particles the $\pcdV$ distribution has a first quartile of \qty{-1.21}{\percent}, a third quartile of \qty{-0.09}{\percent}, and an inter-quartile range (IQR) of \qty{1.12}{\percent}. Applying the standard Tukey $1.5\times$IQR fence of \cref{eq:tukey} to identify outliers yields a lower fence at \qty{-2.89}{\percent} and an upper fence at \qty{+1.59}{\percent}. Six particles fall below the lower fence and are flagged as outliers; no particle falls above the upper fence.

\begin{figure}[htbp]
  \centering
  \includegraphics[width=\textwidth]{fig_volume_sphericity.png}
  \caption[Per-particle volume change versus sphericity change]{Per-particle relative volume change $\pcdV$ (horizontal axis) plotted against per-particle relative sphericity change $\pcdPsi$ (vertical axis), for all 45 tracer particles, with outliers (Mode~B) identified.}
  \label{fig:vol-sph}
\end{figure}

The Mode~A / Mode~B classification used to describe these two populations in earlier documentation is retained as a labelling convenience throughout this chapter---Mode~A refers to the 39 inlier particles (within the Tukey fence) and Mode~B refers to the six outlier particles outside it---but no further mechanistic content is read into the labels themselves. The remainder of the chapter reports the morphological response of the inlier population and treats the outliers as a separately identified set whose contribution to each descriptor is quoted but not averaged into the inlier statistics.

\begin{table}[htbp]
  \centering
  \caption[Per-test summary of the volume-change response]{Per-test summary of the volume-change response. $n_{\mathrm{in}}$ is the number of inlier (Mode~A) particles, $n_{\mathrm{out}}$ the number of outlier (Mode~B) particles; the mean and median $\pcdV$ are computed on the inlier population.}
  \label{tab:volume-change}
  \begin{tabular}{@{}c l c c S[table-format=-1.2] S[table-format=-1.2]@{}}
    \toprule
    \colhead{Test} & \colhead{Configuration} & \colhead{$n_{\mathrm{in}}$} & \colhead{$n_{\mathrm{out}}$} & {\colhead{Mean $\pcdV$}} & {\colhead{Median $\pcdV$}}\\
    \midrule
    \test{T1} & \cfg{PUBM-FB-300} & 5 & 0 & -0.18 & -0.20\\
    \test{T2} & \cfg{RUBM-FB-450} & 5 & 0 & -0.10 & -0.09\\
    \test{T3} & \cfg{UR-FB-300}   & 5 & 0 & -0.44 & -0.44\\
    \test{T4} & \cfg{UR-RB-450}   & 3 & 2 & -0.87 & -1.24\\
    \test{T5} & \cfg{RUBM-RB-450} & 4 & 1 & -0.84 & -0.25\\
    \test{T6} & \cfg{PUBM-RB-450} & 4 & 1 & -0.30 & -0.07\\
    \test{T7} & \cfg{UR-RB-450}   & 4 & 1 & -0.71 & -0.45\\
    \test{T8} & \cfg{RUBM-RB-300} & 4 & 1 & -0.58 & -0.04\\
    \test{T9} & \cfg{PUBM-RB-300} & 5 & 0 & -0.18 & -0.15\\
    \bottomrule
  \end{tabular}
\end{table}

\section{Shape evolution: sphericity, convexity, angularity, and roundness}
\label{sec:shape-evolution}

Angularity and roundness are the two principal descriptors of the sharpness of corner-and-edge features on a particle, while sphericity and convexity integrate over the whole grain. Both classes are expected to evolve toward smoother values under abrasive impact loading, though the corner-sensitive descriptors depend on the subtler local curvature of the surface mesh rather than the global aspect ratio.

\Cref{fig:angularity} plots the angularity index before versus after the test for all 45 particles; the angularity values reported here are negative scalars, so less-negative (closer to zero) values indicate smoother corners. The mechanistic interpretation is that the angularity descriptor is sensitive to local surface curvature and is therefore strongly affected by which corner is preferentially abraded on a given particle. A particle with two sharp corners may lose only one of them under twenty blows; the remaining corner is captured by the maximum-curvature feature in the post-test scan and the descriptor records little change, or even a small increase if the remaining corner becomes more pronounced relative to the now-smoothed neighbour. This within-particle redistribution of curvature is not visible at the level of the global sphericity---which integrates over the whole particle---but is captured directly by angularity. The scatter in \cref{fig:angularity} therefore does not contradict the smoothing-by-abrasion picture established by sphericity; it adds the observation that the smoothing is locally selective, not uniformly distributed across all corners of all particles.

\begin{figure}[htbp]
  \centering
  \includegraphics[width=0.91\textwidth]{fig_angularity.png}
  \caption[Angularity index before versus after]{Angularity index before versus after the impact-loading test, for all 45 tracer particles. The angularity values are negative scalars; less-negative (closer to zero) values indicate smoother corners.}
  \label{fig:angularity}
\end{figure}

The sphericity---the classical Wadell descriptor \citep{wadell1932}---is the single most informative shape descriptor for diagnosing corner abrasion. A particle that loses its sharp corners without otherwise changing shape becomes more spherical. \Cref{fig:sphericity} plots the paired sphericity---$\Psi_{\mathrm{after}}$ on the vertical axis against $\Psi_{\mathrm{before}}$ on the horizontal axis---for all 45 tracer particles. The unit-slope line $\Psi_{\mathrm{after}} = \Psi_{\mathrm{before}}$ is drawn for reference; points above this line indicate a sphericity increase, points on it indicate no change, and points below it a sphericity decrease.

\begin{figure}[htbp]
  \centering
  \includegraphics[width=0.91\textwidth]{fig_sphericity.png}
  \caption[Sphericity before versus after]{Sphericity before versus after the impact-loading test, for all 45 tracer particles.}
  \label{fig:sphericity}
\end{figure}

The systematic displacement of the points above the unit-slope line is the single most consistent finding of this chapter. The mean paired sphericity change across all 45 particles is $\dPsi = +0.0100$, and the paired $t$-test against the null of zero mean change yields $t = 3.23$, $p = 0.002$. Restricting attention to the 39 inlier particles only, the mean change is $\dPsi = +0.0041$ and the $t$-statistic is $6.05$ with $p < 10^{-6}$. The smaller mean for the inliers reflects the fact that the outlier particles, after their fragmentation, are individually more spherical than they were before; their removal slightly reduces the average rise but, more importantly, sharpens the significance of the rise considerably because the within-population variance falls. The conclusion is that sphericity rises in essentially every particle regardless of configuration. Corner abrasion is therefore a universal response to the drop-weight impact protocol, and no configuration in the matrix---including the best-performing \test{T1} (\cfg{PUBM-FB-300})---eliminates it. This is consistent with the abrasion-driven rounding observed in cyclic loading by \citet{qian2014} and \citet{guo2018}, and in drop-weight impact loading by \citet{koohmishi2020}.

\Cref{fig:convexity} reports the convexity $C = V/\Vch$ before versus after the test, with the same layout and statistics as for the sphericity. Convexity rises in parallel with sphericity, confirming that the particles become smoother and more compact as their corners are abraded.

\begin{figure}[htbp]
  \centering
  \includegraphics[width=0.91\textwidth]{fig_convexity.png}
  \caption[Convexity before versus after]{Convexity $C = V/\Vch$ before versus after the impact-loading test, for all 45 tracer particles. Layout and statistics as for \cref{fig:sphericity}.}
  \label{fig:convexity}
\end{figure}

The roundness response, by contrast, is mixed (\cref{fig:roundness}). The mean paired change is $+0.038$ across all 45 particles ($t = 0.67$, $p = 0.51$), and $+0.058$ across the 39 inliers ($t = 1.00$, $p = 0.32$). Like angularity, roundness is sensitive to local corner curvature and reflects the selective nature of corner abrasion: most particles experience either a modest increase or a modest decrease in roundness, depending on which corner the impact-loading process abraded most. The sphericity-and-convexity rise reported above is the more reliable population-level diagnostic of the abrasion process; the angularity-and-roundness scatter reports the corresponding particle-level heterogeneity.

\begin{figure}[htbp]
  \centering
  \includegraphics[width=0.91\textwidth]{fig_roundness.png}
  \caption[Roundness before versus after]{Roundness before versus after the impact-loading test, for all 45 tracer particles.}
  \label{fig:roundness}
\end{figure}

\section{Discussion}
\label{sec:morph-discussion}

The morphological script computes twenty-two descriptors per mesh, of which the four reported in \cref{sec:shape-evolution} are the principal shape descriptors. To provide a global view of the pattern of mean change across all configurations and a wider set of descriptors, \cref{fig:heatmaps} collects the mean per-particle percentage change of fourteen descriptors across the nine tests in two side-by-side heat maps. Outlier particles are removed before the means are computed; the visualisation is therefore a Mode~A (inlier) picture of the morphological response. The upper panel shows the size and primary shape descriptors, which have a comparatively narrow dynamic range (less than $\pm\qty{2.5}{\percent}$); the lower panel shows the roundness and shape-ratio descriptors, which span a wider range (up to about $\pm\qty{80}{\percent}$).

\begin{figure}[htbp]
  \centering
  \includegraphics[width=\textwidth]{fig_heatmaps.png}
  \caption[Heat maps of mean per-particle change across descriptors]{Heat maps of mean per-particle percentage change in fourteen morphological descriptors across the nine tests, with outlier particles removed. Upper panel: size and primary shape descriptors. Lower panel: roundness and shape-ratio descriptors.}
  \label{fig:heatmaps}
\end{figure}

Three patterns are visible in the upper panel of \cref{fig:heatmaps}. First, all the size descriptors (volume, surface area, convex-hull volume, convex-hull surface, equivalent diameter) are uniformly negative across the nine configurations. The sign-uniformity is mechanistically reassuring: inlier particles do not gain mass under impact loading. The magnitude of the size loss tracks the ranking of \cref{tab:volume-change}, with the rigid-base \qty{450}{\mm} tests \test{T4}, \test{T5}, and \test{T7} showing the largest size losses. Second, sphericity and convexity are uniformly positive across the nine configurations, consistent with the universal abrasion-driven rounding established in \cref{sec:shape-evolution}. Third, the surface-texture-index column shows small but mixed-sign changes (between \qty{-0.42}{\percent} and \qty{+0.17}{\percent}), with no clear configuration-level pattern; this descriptor is sensitive to the highest-frequency mesh asperities and is dominated by scan-to-scan reconstruction noise rather than the impact-loading signal.

The lower panel of \cref{fig:heatmaps} reports the roundness and shape-ratio descriptors, which exhibit a much larger dynamic range. The angularity column in particular spans from \qty{-25}{\percent} at \test{T5} to \qty{+79}{\percent} at \test{T9}---a span of more than 100 percentage points across the nine configurations. The sign of the angularity change does not group obviously by base condition or by drop height: \test{T1}, \test{T4}, \test{T6}, \test{T8}, and \test{T9} are all positive (sharper after the test, in the convention adopted here), while \test{T2}, \test{T3}, \test{T5}, and \test{T7} are negative (smoother). The roundness column shows the same mixed-sign pattern, with the largest positive value at \test{T9} (\qty{+27}{\percent}) and the largest negative value at \test{T5} (\qty{-13}{\percent}). The aspect-ratio and flatness columns are similarly heterogeneous. The mixed-sign pattern of these descriptors confirms the observation made in \cref{sec:shape-evolution} about angularity: corner abrasion is distributed selectively rather than uniformly across the corners of each particle, and the population-level mean of these locally sensitive descriptors is therefore not a reliable diagnostic of the configuration-level damage. The size descriptors and the global shape descriptors (sphericity and convexity) of the upper panel are the more robust diagnostics.
